\chapter{SGLang}
SGLANG (\cite{zheng2024sglangefficientexecutionstructured}) is a system for efficient execution of complex language model programs. SGLang consists of a frontend language and a runtime. The frontend simplifies programming with primitives for generation and parallelism control. The runtime accelerates execution with novel optimizations like RadixAttention for KV cache reuse and compressed finite state machines for faster structured output decoding.

\section{LM Programs}
The emergence of few-shot learning, self-consistency, skeleton-of-thought, tree-of-thought and various other patterns signifies a shift in our interaction with LLMs, moving from simple chatting to a more sophisticated form of programmatic usage of LLMs, which means using a program to schedule and control the generation processes of LLMs. These programs are referred to as LM (Language Model) programs. There are two common properties of LM programs: 
\begin{itemize}
    \item LM programs typically contain multiple LLM calls interspersed with control flow. This is needed to complete complex tasks and improve overall quality.
    \item LM programs receive structured inputs and produce structured outputs. This is needed to enable the composition of LM programs and to integrate LM programs into existing software systems.
\end{itemize}


\section{Programming Model}

\textbf{Language primitives:} SGLang provides primitives for controlling prompt state, generation, and parallelism. Here are the primitives: ``gen'' calls a model to generate and stores the results in a variable with the name specified in its first argument. It supports a ``regex'' argument to constrain the output to follow a grammar defined by a regular expression (e.g., a JSON schema). ``select'' calls a model to choose the highest probability option from a list.  The operator ``+='' or  ``extend'' appends a string to the prompt. The operator ``[variable\_name]'' fetches the results of a generation. ``fork'' creates parallel forks of the prompt state. ``join'' rejoins the prompt state. ``image'' and ``video'' take in image and video inputs.

\begin{figure}[H]
    \centering
    \includegraphics[width=0.8\textwidth]{\FIGDIR/sglang_model.png}
    \caption{SGLANG Frontend Language Model (\cite{zheng2024sglangefficientexecutionstructured})}
\end{figure}



\section{RadixAttention}
A radix tree is a data structure that serves as a space-efficient alternative to a classical trie (prefix tree). Unlike typical trees, the edges of a radix tree can be labeled not just with single elements but also with sequences of elements of varying lengths, significantly enhancing efficiency. In SGLANG, a radix tree is utilized to manage a mapping between sequences of tokens, and their corresponding KV cache tensors. These KV cache tensors are stored in a non-contiguous, paged layout, where the size of each page is equivalent to one token. Because GPU memory is quickly filled by the KV cahce, a simple LRU eviction policy is introduced that evicts the least recently used leaf first. By evicting leaves first, we enable the re-use of their common ancestors until those ancestors become leaves and are also evicted. 

In the continuous batching setting, we cannot evict nodes used by the currently running batch. Each node maintains a reference counter indicating how many running requests are using it. A node is evictable if its reference counter is zero.

The cache hit rate is defined as $\frac{\texttt{number of cached prompt tokens}}{\texttt{number of prompt tokens}}$. If the request scheduler frequently switches between different, unrelated requests, it can lead to cache thrashing and a low hit rate.

In the batch-processing setting the requests are sorted by matched prefix length and prioritize requests with longer matched prefixes instead of using a first-come, first-served schedule. 

\begin{figure}[H]
    \centering
    \includegraphics[width=\textwidth]{\FIGDIR/sglang_radix_tree.png}
    \caption{SGLANG Radix Tree Example (\cite{zheng2024sglangefficientexecutionstructured})}
\end{figure}


\section{Compressed Finite State Machine}
In LM programs, users often want to constrain the model's output to follow specific formats, such as JSON schemas. Existing systems support this by converting a regular expression into a finite state machine (FSM). During decoding, they maintain the current FSM state, retrieve allowed tokens from the next states, and set the probability of invalid tokens to zero, decoding token by token. They can only decode one token at a time because the lack of integration between the FSM and the model runner in existing systems prevents multi-token processing, resulting in slow decoding. SGLang overcomes this limitation by creating a fast constrained decoding runtime with a compressed FSM. This runtime analyzes the FSM and compresses adjacent singular-transition edges in the FSM into single edges.


\section{FastTree: Optimizing Radix Tree based KV Cache Computation}
\label{sec:fasttree}
Existing systems still use conventional attention kernels that concatenate the KV cache per query and process each query independently. This leads to redundant memory loads of shared contexts from slow global memory and poor tensor core utilization due to mismatched tile dimensions. 

\begin{wrapfigure}{r}{0.5\textwidth}
    \centering
    \vspace{-10pt}
    \includegraphics[width=\linewidth]{\FIGDIR/kv_cache_sharing_through_radix_tree.png}
    \caption{KV cache sharing through Radix Tree (\cite{pan2025fasttree})}
\end{wrapfigure}

FastTree (\cite{pan2025fasttree}) introduces GPU kernels specifically designed to process queries that share common prefixes within a radix tree structure. By adapting the computation to the tree topology at runtime, FastTree eliminates redundant memory operations and improves throughput.

\textbf{Tree-structure Adaptive Runtime Optimization.}
Different ways of grouping queries—based on how prefixes are merged—can lead to drastically different computation patterns and memory-sharing behavior. Even with identical tree topologies, certain grouping plans might waste computation due to tile padding or cause additional global memory transactions when intermediate results must be written out. FastTree formulates this as an optimization problem: given a KV cache tree $T=(V,E)$, it searches for an edge assignment function $f:E\rightarrow\{0,1\}$ that minimizes the end-to-end latency of the GPU attention kernel, balancing the trade-off between computation reuse and intermediate data cost.  

Padding overhead arises when the number of queries or context tokens fails to meet the tensor core tile size, wasting both shared memory and FLOPs. This is modeled through cost functions $C_{P,q}$ and $C_{P,c}$, which estimate padding-induced inefficiency at the query and context dimensions respectively. Similarly, splitting the KV matrices introduces intermediate reduction overheads, modeled as $\text{SplitKV\_Cost}$, capturing both padding and extra memory I/O cost. FastTree employs a greedy runtime algorithm (Fig.~\ref{fig:fasttree_alg}) that evaluates these costs per grouping plan and dynamically selects the most efficient configuration.

\begin{wrapfigure}{r}{0.5\textwidth}
    \centering
    \vspace{-10pt}
    \includegraphics[width=\linewidth]{\FIGDIR/fast_tree_algorithm.png}
    \caption{Greedy runtime optimization algorithm (\cite{pan2025fasttree})}
    \label{fig:fasttree_alg}
\end{wrapfigure}

\textbf{Efficient Attention Kernel Design.}
After determining the optimal grouping plan, FastTree performs the attention computation tile-by-tile within each group. It parallelizes the tiles of the query (Q) matrix across thread blocks, while iterating over key (K) and value (V) tiles sequentially within each block due to inter-tile dependencies. By caching QKV tiles in shared memory and applying an online softmax, FastTree minimizes global memory transactions. Furthermore, aggregating multiple queries enables the transformation of GEMV into GEMM operations, thereby unlocking high tensor core utilization on modern GPUs.

Overall, FastTree demonstrates that by coupling radix tree–based KV cache organization with adaptive runtime planning and GPU-aware attention kernels, significant reductions in redundant computation and memory access can be achieved, leading to markedly higher throughput and lower latency for shared-context inference workloads.
